\documentclass[12pt]{amsart}
\usepackage[all]{xy} % for commutative diagrams
\usepackage{amsmath} % various goodies
\usepackage{amsthm} % some stuff for theorem numbering
\usepackage{amsfonts}

%% Theorem-type environments.
% The first parameter to \newtheorem is what the environment is called in the 
% LaTeX source. (What goes into the braces after \begin and \end.)
% The second is the name that shows up in the text.

\theoremstyle{plain} 
\newtheorem{theorem}{Theorem}
\newtheorem{lemma}{Lemma}
\newtheorem{corollary}{Corollary}
\newtheorem{proposition}{Proposition}
\newtheorem{conjecture}{Conjecture}
% (Criterion, Algorithm)

\theoremstyle{definition} 
\newtheorem{definition}{Definition}
\newtheorem{example}{Example}
% (Condition, Problem)

\theoremstyle{remark} 
\newtheorem{remark}{Remark}
\newtheorem{note}{Note}
% (Notation, Claim, Summary, Acknowledgement, Case, Conclusion)

%% Math operators (easy thanks to amsmath)
\DeclareMathOperator{\Aut}{Aut}
\DeclareMathOperator{\core}{Core}
\DeclareMathOperator{\orb}{Orb}
\DeclareMathOperator{\cay}{Cay}
\DeclareMathOperator{\sab}{Sab}
\DeclareMathOperator{\fib}{fiber}



% Commands:
\newcommand{\cB}{\mathcal{B}}
\newcommand{\cS}{\mathcal{S}}
\newcommand{\cP}{\mathcal{P}}
\newcommand{\cZ}{\mathcal{Z}}
\newcommand{\cD}{\mathcal{D}}
\newcommand{\cH}{\mathcal{H}}
\newcommand{\cQ}{\mathcal{Q}}
\newcommand{\cK}{\mathcal{K}}
\newcommand{\cN}{\mathcal{N}}

% Lazy:
\newcommand{\travH}{\{ \tau_1, \tau_2, \dots \tau_n \}}
\newcommand{\travR}{\{ t_1, t_2, \dots t_b \}}
\newcommand{\rangen}{\{1 \dots n \}}

\begin{document}

\section{Description of the problem}
We would like to understand the ways in which a symmetric graph
$(G, \Gamma)$ can be ``unfolded'' into a larger imprimitive symmetric graph
$(\widetilde{G}, \widetilde{\Gamma})$ which admits the original graph as a quotient.
On a group theoretic level, if we knew the subgroup structure of $G$
and we knew all the groups $\widetilde{G}$ which have $G$ as a
composition factor then in a sense we would ``know'' all the possible unfoldings.
What would still be lacking is a {\it combinatorial} understanding
of how the structure of the new graph is related to that of the original.

The {\it three-arc graph construction} gives a nice description
of the graph $\widetilde{\Gamma}$ in terms of the graph $\Gamma$
in the special case where $\Gamma \cong \sab(G,H,a)$,
$\widetilde{\Gamma} \cong \sab(G,K,a)$ and $K \cong a^{-1}Ha \cap H$.
That is, the special case where the stabilizer of an arc of $\Gamma$
is isomorphic to the stabilizer of a vertex of $\widetilde{\Gamma}$.
One may ask whether there other {\it group theoretic} assumptions
which may be imposed on a pair of symmetric graphs $(\Gamma, \widetilde{\Gamma})$
which will yield a nice description of $\widetilde{\Gamma}$ in terms of $\Gamma$.

I have considered the case where $\Gamma \cong \sab(G,H,a)$,
$\widetilde{\Gamma} \cong \sab(G,K,a)$ 
and there exists a normal subgroup $N$ of $G$ such that
$G$ is a semidirect product of $N$ by $H$.
Here are some rough notes.

\section{Labelling technique}

Every $g \in G$ has a unique expression of the form $g = (h,n)$
where $h \in H$ and $n \in N$.
Since the elements of $N$ form a traversal of the cosets of $H$ in $G$,
the vertices of $\Gamma$ may labelled by elements of $N$
and in fact form a group.
The action of $G$ on the vertices of $\Gamma$ is given by
$m^{(h,n)} = h^{-1}mhn$.
Let us take $\cD$ to be the the design induced on the fiber of the
vertex of $\Gamma$ which is labelled by the identity of $N$.

The points of $\widetilde{\Gamma}$ may be labelled by ordered pairs of the form
$(x, m)$ where $x$ is a coset of $K$ and $m \in N$.
The action of $G$ on the points of $\widetilde{\Gamma}$ is then given by: 
$(x,m)^{(h,n)} = (x^h, h^{-1}mhn)$. 
This is indeed an action, since:
\begin{eqnarray*}
{(x,m)^{(h_1,n_1)}}^{(h_2,n_2)} & = & (x^{h_1}, h_{1}^{-1} m h_1 n_1)^{(h_2,n_2)} \\
                                & = & (x^{h_1 h_2}, h_{2}^{-1}h_1^{-1} m h_1 n_1 h_2 n_2) \\
                      & = & (x^{h_1 h_2}, h_{2}^{-1}h_1^{-1} m h_1 h_2 h_2^{-1} n_1 h_2 n_2) \\
		      & = & (x,m)^{(h_1 h_2, h_2^{-1} n_1 h_2 n_2)} \\
		      & = & (x,m)^{(h_1,n_1)(h_2,n_2)} \\
\end{eqnarray*}


\section{Self-paired orbital on the flags $\cD$}

I will show that there is a self-paired orbital on the flags
of $\cD$ which can be used to reconstruct $\widetilde{\Gamma}$
from $\Gamma$.

Suppose that $((x,p), (y,b))$ is an arc $\widetilde{\Gamma}$.
Then $(p,q)$ is an arc of the quotient $\Gamma$.
Making use of the automorphism $(1,p^{-1})$ of $\widetilde{\Gamma}$
we have that $((x,p), (y,b))^{(1,p^{-1})} = ((x,1), (y,bp^{-1}))$ 
is also an arc of $\widetilde{\Gamma}$,
and so $(x,bp^{-1})$ is a flag of $\cD$

Since $\widetilde{\Gamma}$ is assumed to be a simple graph,
if $((x,p), (y,b))$ is an arc of $\widetilde{\Gamma}$
then so is $((y,b),(x,p))$.
Making use of the automorphism $(1,b^{-1})$ we have that
and $((y,b),(x,p))^{(1,b^{-1})} = ((y,1), (x,pb^{-1}))$
is an arc of $\widetilde{\Gamma}$.
and so $(y, pb^{-1})$ is a flag of $\cD$.

What we have done is ``decomposed'' the arc $((x,p), (y,b))$
of $\widetilde{\Gamma}$ into the arc $(p,q)$ of $\Gamma$
together with the pair of flags $((x,bp^{-1}),(y, pb^{-1}))$
of $\cD$.
Let $\Delta$ be the {\it orbital} on the flags of $\cD$
which contains the pair $((x,bp^{-1}),(y, pb^{-1}))$.
That is $\Delta = \{ ((x,bp^{-1}),(y, pb^{-1}))^h : h \in H \}$
\hfill \break

Suppose we choose a different arc $((w,q),(z,d))$ of $\widetilde{\Gamma}$.
Then by arguments identical to those above $(q,d)$ is an arc of $\Gamma$
and $(w,dq^{-1})$ and $(z,qd^{-1})$ are flags of $\cD$.
Since $\widetilde{\Gamma}$ is $G$-arc transitive, we can find some $g = (h,n)$
such that $((x,p),(y,b))^{(h,n)} = ((w,q),(z,d))$.
That is:
\begin{eqnarray*}
w & = & x^h \\
z & = & y^h \\
q & = &  h^{-1}phn \\
d & = &  h^{-1}bhn \\
\end{eqnarray*}

So: 
\begin{eqnarray*}
(w,dq^{-1}) & = & (x^h, h^{-1}bhnn^{-1}h^{-1}p^{-1}h) \\
            & = & (x^h, h^{-1}bp^{-1}h) \\
	    & = & (x,bp^{-1})^h.
\end{eqnarray*}
Simmilarly:
\begin{eqnarray*}
(z,qd^{-1}) & = & (y^h, h^{-1}phnn^{-1}h^{-1}b^{-1}h) \\
            & = & (y^h, h^{-1}pb^{-1}h) \\
	    & = & (y,pb^{-1})^h.
\end{eqnarray*}

\newpage
It follows immediately that $((w,dq^{-1}),(z,qd^{-1})) \in \Delta$.
Thus, any arc of $\widetilde{\Gamma}$ can be ``decomposed''
into an arc of $\Gamma$ together with a pair of flags of $\cD$
contained in the orbital $\Delta$.
\hfill \break

I will now show that this orbital $\Delta$ on the flags of $\cD$
is self-paired.
Since $\widetilde{\Gamma}$ is simple.
There must be some $g = (h,n)$ which ``flips'' the arc
$((x,p), (y,b))$. That is, there must be some $g = (h,n)$
such that $((x,p), (y,b))^{(h,n)} = ((y,b), (x,p))$.
By an argument identical to the one on the previous page, 
we must have that

\begin{eqnarray*}
(y,pb^{-1}) & = & (x,bp^{-1})^h \\
(x,bp^{-1}) & = & (y,pb^{-1})^h
\end{eqnarray*}

\hfill \break
This shows that the orbital $\Delta$ is self-paired.

\section{Reconstruction}
Given $\Gamma$ and $\cD$ together with a self paired orbital $\Delta$
on the flags of $\cD$ it should be possible to reconstruct $\widetilde{\Gamma}$.
Let $N(1)$ denote the neighbourhood of the vertex ``$1$'' of the $\Gamma$
and let $P$ and $B$ denote the points and blocks respectively 
of the design $\cD$.
Let $\eta : N(1) \to B$ be some map establishing a permutation isomorphism
between the permutation groups $(H,N(1))$ and $(H,B)$.

The graph $\widetilde{\Gamma}$ can be described as the graph
with vertex set $P \times V \Gamma$
and arc set $((x,v), (y,w))$ such that 
\begin{enumerate}
\item $(v,w)$ is an arc of $\Gamma$
\item $(x,\eta(vw^{-1})$ is a flag of $\cD$
\item $(y, \eta(wv^{-1})$ is a flag of $\cD$
\item $((x,\eta(vw^{-1}),(y, \eta(wv^{-1})) \in \Delta$ 
\end{enumerate}
\hfill \break

The action of $G$ on $P \times V \Gamma$ is that given in section 2.
Since $\cD$ is $H$-flag transitive,
for any pair of vertices $(x,p)$ and $(y,q)$,
we can find some $h \in H$ such that $x^h = y$. 
Now: 
\begin{eqnarray*}
(x,p)^{(h,h^{-1}p^{-1}hq)} & = & (x^h, h^{-1}phh^{-1}p^{-1}hq) \\
                           & = & (y,q).
\end{eqnarray*}
So this action is transitive.
\hfill \break

We need to check that this action is well defined on the arcs.
Suppose that $((x,p),(y,q))$ is any arc and $(h,n)$ is 
any element of $G$.
We need to check that 
$((x,p),(y,q))^{(h,n)}$ is also an arc.

Since $((x,p),(y,q))$ is an arc of $\widetilde{\Gamma}$
we know that $(p,q)$ is an arc of $\Gamma$
and that: 
$$((x,\eta(pq^{-1}), (y,\eta(qp^{-1})) \in \Delta$$
\hfill \break

The arc $((x,p),(y,q))^{(h,n)}$ of $\widetilde{\Gamma}$
decomposes into the arc $(p,q)^{(h,n)}$ of $\Gamma$.
and the pair of flags: $$((x,\eta(pq^{-1})^h, (y,\eta(qp^{-1}))^h \in \Delta$$
So the action is well-defined on arcs.
\hfill \break

Finally we need to check that this action is transitive on the arcs.
If $((x,p),(y,b))$ and $((w,q),(z,d))$ are any two arcs of $\widetilde{\Gamma}$
then we must have: 

$$((x,\eta(pb^{-1}),(y,\eta(bp^{-1})) \in \Delta$$

and:  

$$((w,\eta(qd^{-1}),(z,\eta(dq^{-1})) \in \Delta.$$

\hfill \break
So there must be some $h \in H$ such that:
$$((x,\eta(pb^{-1}),(y,\eta(bp^{-1}))^h = ((w,\eta(qd^{-1}),(z,\eta(dq^{-1})).$$

\hfill \break
Since $\eta$ is a permutation homomorphism,
$\eta(bp^{-1})^h = \eta(dq^{-1})$ implies that
$\eta(h^{-1}bp^{-1}h) = \eta(dq)$ which implies that
$h^{-1}bp^{-1}h = dq^{-1}$.
Now we have: 
\begin{eqnarray*}
(x,p)^{(h,h^{-1}p^{-1}hq)} & = & (x^h, h^{-1}phh^{-1}p^{-1}hq) \\
                           & = & (w,q)
\end{eqnarray*}
And: 
\begin{eqnarray*}
(y,b)^{(h,h^{-1}p^{-1}hq)} & = & (y^h, h^{-1}bhh^{-1}p^{-1}hq) \\
                           & = & (z, h^{-1}bp^{-1}hq) \\
			   & = & (z, dq^{-1}q) \\
			   & = & (z,d) \\
\end{eqnarray*}
Thus the action is transitive on the arcs as claimed.

\newpage
\section{Another Special Case}

Suppose that $\Gamma = \sab(G,H,a)$.
Let $\overline{H} = a^{-1}Ha \cap H$,
so the local permutation group induced at each vertex of $\Gamma$
is equivalent to $(H, \cos(\overline{H}))$.
Suppose further that there is some $K$ such that $\overline{H} < K < H$
and $a \not\in K$. Then the graph $\widetilde{\Gamma} = \sab(G,K,a)$
is an extension of $\Gamma$.
Let $\overline{K} = a^{-1}Ka \cap K$.
Since $K < H$ we must also have $a^{-1}Ka < a^{-1}Ha$
so $\overline{K} \leq \overline{H}$.
Since $\overline{H} < K$ we must also have
$a^{-1}\overline{H}a < a^{-1}Ka$
but since $a$ is an involution $a^{-1}\overline{H}a = \overline{H}$
so we have $\overline{H} \leq \overline{K}$.
That is $\overline{H} = \overline{K}$
and the local permutation group induced at each vertex of
$\widetilde{\Gamma}$ is $(K, \cos(\overline{H}))$.

The pair $(\Gamma,\widetilde{\Gamma})$ satisfy the property that
{\it globally} $\widetilde{\Gamma}$ is imprimitive,
admitting $\Gamma$ as quotient, but {\it locally}
$\Gamma$ is imprimitive, and the local permutation group of
$\widetilde{\Gamma}$ is a quotient of the local permutation
group of $\Gamma$. We wish to understand how $\Gamma$
and $\widetilde{\Gamma}$ are related {\it structurally},
and ideally find some method of constructing $\widetilde{\Gamma}$
from $\Gamma$.

As a preliminary observation, let $n = [G:H]$ be the number of vertices of $\Gamma$,
let $v = [K:\overline{H}]$ be the valency of $\widetilde{\Gamma}$
and let $r = [H:K]$.
The graph $\widetilde{\Gamma}$ has $[G:K] = [G:H][H:K] = nr$
vertices. That is, $r$ times as many vertices as $\Gamma$.
The valency of $\Gamma$ is $[H:\overline{H}] = [H:K][K:\overline{H}] = rv$.
That is $r$ times the valency of $\widetilde{\Gamma}$.
Since the number of edges in a graph is equal to half the
valency times the number of vertices, it follows that
both $\Gamma$ and $\widetilde{\Gamma}$ have the same number of edges.
\hfill \break

The action of $G$ on the arcs of $\Gamma$ is permutation equivalent
to $(G, \cos(\overline{H}))$.
Since $\Gamma$ is simple, each arc $(\alpha,\beta)$ has a {\it pair}
namely $(\beta,\alpha)$.
The function $\varphi : A \Gamma \to A \Gamma$ 
which sends each arc to its pair is an involution
It also preserves the action of $G$.

Now, since $\overline{H} < H < G$ the action of $G$
on the arcs of $\Gamma$ is imprimitive.
That is, there is some $G$-invariant partition $\cB$
of $A \Gamma$ such that the action of $G$ on $\cB$
is permutation equivalent to $(G, \cos(H))$.
In fact, this is the partition $\cB = \{ B(\alpha) : \alpha \in V \Gamma \}$
where $B(\alpha)$ is the set of all arcs who's initial vertex is $\alpha$.

Let $\pi : A \Gamma \to \cB$ be the map which sends
each arc to the block of $\cB$ containing it.
We may define a new graph $\Sigma$ we vertex set
$\cB$ where $p$ is adjacent to $b$ if and only if
there is some $x \in \pi^{-1}(p)$ and some $y \in \pi^{-1}(b)$
such that $y = \varphi(x)$.
This new graph $\Sigma$ is in fact isomorphic to $\Gamma$.

The action of $G$ on the arcs of $\Gamma$ is permutation
isomorphic to the action of $G$ on the arcs of $\widetilde{\Gamma}$.
Both actions are of the form $(G, \cos(\overline{H}))$
Since $\overline{H} < K < G$ there exists some $G$-invariant partition
$\mathfrak{B}$ of $A \Gamma$ such that the action of $G$ on $\mathfrak{B}$
is permutation equivalent to $(G, \cos(H))$.
The partition $\mathfrak{B}$ is a refinement of the partition $\cB$.

Let $\lambda : A \Gamma \to \mathfrak{B}$ be the map hich sends
each arc to the block of $\mathfrak{B}$ containing it.
We may define a new graph $\widetilde{\Sigma}$ we vertex set
$\mathfrak{B}$ where $p$ is adjacent to $b$ if and only if
there is some $x \in \lambda^{-1}(p)$ and some $y \in \lambda^{-1}(b)$
such that $y = \varphi(x)$.
This new graph $\Sigma$ is in fact isomorphic to $\widetilde{\Gamma}$.
Thus we have a method of construction $\widetilde{\Gamma}$ from $\Gamma$.

\end{document}
